% ====================================================================
% VİZE RAPORU – Türk Masalları Derlemesi Web Sayfası
% DERLEME: Overleaf'te Menu > Compiler > XeLaTeX (pdfLaTeX de çalışır).
% ====================================================================
\documentclass[12pt,a4paper]{article}

\usepackage{iftex}
\ifPDFTeX
  \usepackage[T1]{fontenc}
  \usepackage[utf8]{inputenc}
\else
  \usepackage{fontspec}
\fi
\usepackage[turkish]{babel}
\usepackage[margin=2.4cm]{geometry}
\usepackage{setspace}
\usepackage{graphicx}
\usepackage{booktabs}
\usepackage{longtable}
\usepackage{array}
\usepackage{colortbl}
\usepackage{xcolor}
\usepackage{enumitem}
\usepackage{listings}
\usepackage{amssymb}
\usepackage{caption}
\usepackage{fancyhdr}
\usepackage{hyperref}

\onehalfspacing
\setlength{\parindent}{0pt}
\setlength{\parskip}{6pt}
\setlist{itemsep=2pt, topsep=4pt}
\hypersetup{colorlinks=true, linkcolor=blue!50!black, urlcolor=blue!50!black}

\definecolor{ana}{HTML}{14325C}
\definecolor{acik}{HTML}{EEF2F8}

\lstset{basicstyle=\ttfamily\small, backgroundcolor=\color{acik},
  frame=single, rulecolor=\color{ana}, breaklines=true,
  columns=fullflexible, keepspaces=true, showstringspaces=false}

% Tablo ayarları
\renewcommand{\arraystretch}{1.3}

% Üst ve alt bilgi (Header/Footer)
\pagestyle{fancy}
\fancyhf{}
\fancyhead[L]{\textcolor{gray}{\small Türk Masalları}}
\fancyhead[R]{\textcolor{gray}{\small Vize Raporu}}
\fancyfoot[C]{\thepage}
\renewcommand{\headrulewidth}{0.4pt}
\renewcommand{\footrulewidth}{0pt}

% ---------- Yardımcı komutlar ----------
\newcommand{\aciklama}[1]{{\color{gray}\itshape (#1)}}

% Haftalık tablo görseli: dosya yoksa "görsel ekleyin" kutusu çıkar
\newcommand{\haftagorsel}[1]{%
  \IfFileExists{#1}{%
    \includegraphics[width=\linewidth,height=4.2cm,keepaspectratio]{#1}%
  }{%
    \fbox{\parbox[c][2.6cm][c]{\dimexpr\linewidth-2\fboxsep-2\fboxrule\relax}{%
      \centering\small\color{gray} Ekran görüntüsü / görsel ekleyin\\ \texttt{#1}}}%
  }%
}

% Bölüm 3.4 ekran görüntüsü: dosya varsa ince çerçeveli görsel, yoksa uyarı kutusu
% #1 = dosya adı, #2 = genişlik
\newcommand{\ekrangorsel}[2]{%
  \IfFileExists{#1}{%
    {\setlength{\fboxsep}{0pt}\setlength{\fboxrule}{0.5pt}%
     \fcolorbox{gray!60}{white}{\includegraphics[width=\dimexpr#2-1pt\relax]{#1}}}%
  }{%
    \fbox{\parbox[c][4cm][c]{\dimexpr#2-2\fboxsep-2\fboxrule\relax}{%
      \centering\small\color{gray} Ekran görüntüsü ekleyin\\ \texttt{#1}}}%
  }%
}

% Haftalık katılım kutuları
\newcommand{\katilimkutulari}{\par\vspace{3pt}\small $\square$ Ders içinde katılım sergiledim\par $\boxtimes$ Telafi yaptım}
\newcommand{\katilimkutularibos}{\par\vspace{3pt}\small $\square$ Ders içinde katılım sergiledim\par $\square$ Telafi yaptım}

% Ek görsel komutunu şablonu temiz tutmak için devre dışı bıraktık
\newcommand{\haftaek}[1]{}

% ---------- Kapak bilgileri ----------
\newcommand{\universite}{İstanbul Medipol Üniversitesi / Görsel İletişim Tasarımı}
\newcommand{\dersadi}{Programlamaya Giriş}
\newcommand{\hocaadi}{Meral Şentürk}
\newcommand{\ogrenciadi}{Zeynep Oktav}
\newcommand{\ogrencino}{86230023}
\newcommand{\projeadi}{Türk Masalları}
% Bağlantıları tam adres olarak yazın (https:// ile başlayarak)
\newcommand{\githublink}{\url{https://github.com/oktavzeynep0-lab/turk-masallari-vize}}
\newcommand{\youtubelink}{\url{https://youtu.be/-b3_Y1uaCPI}}

\begin{document}

% ====================================================================
% KAPAK
% ====================================================================
\begin{titlepage}
    \centering
    \vspace*{2cm}

    {\large İstanbul Medipol Üniversitesi / Görsel İletişim Tasarımı}\\[2cm]

    {\Huge \bfseries \color{ana} Türk Masalları}\\[1cm]

    {\Large Türk Masalları Derlemesi Web Sayfası}\\[0.5cm]
    {\large Vize Raporu}\\[3cm]

    {\textbf{Ders:} Programlamaya Giriş}\\[0.3cm]
    {\textbf{Öğretim Üyesi:} Meral Şentürk}\\[1.5cm]

    {\textbf{Hazırlayan:} Zeynep Oktav}\\[0.3cm]
    {\textbf{Öğrenci No:} 86230023}\\[2cm]

    \vfill

    % Proje Bağlantıları Kutusu
    \fbox{\parbox{0.88\linewidth}{\centering\vspace{6pt}
{\large\bfseries Proje Bağlantıları}\\[8pt]
\textbf{GitHub deposu:}\\ \githublink \\[8pt]
\textbf{YouTube sunum videosu:}\\ \youtubelink
\vspace{6pt}}}

    \vfill
    {\large \today \par}
\end{titlepage}

\tableofcontents
\newpage

% ====================================================================
\section{Giriş}
% ====================================================================

\subsection{Proje Özeti}
Bu proje, anonim Türk halk anlatılarından seçilen sekiz masalı tek bir web sayfasında
bir araya getiren dijital bir masal kütüphanesidir. Sayfada her masal, renkli bir kapağı
olan üç boyutlu bir kitap olarak gösterilir; kitaba tıklandığında kapak açılır ve masal
sade bir okuma ekranında sunulur. Kullanıcı masalları başlığa göre arayabilir, türe göre
(hayvan, eğitici, fantastik) süzebilir ve okurken yazı boyutunu değiştirebilir.
Sayfa; HTML, içsel CSS ve içsel JavaScript ile tek bir dosyada hazırlanmıştır.
Masal metinlerinin kurgulanmasında Türkiye Kültür Portalı arşivlerinden yararlanılmıştır.
Bu rapor; masalların seçimini, tasarım kararlarını, kodun işleyişini ve haftalık katkıları
içermektedir.

\subsection{Gereksinimlerin Karşılanması}
\begin{center}
\rowcolors{2}{acik}{white}
\begin{tabular}{@{}p{8.2cm}c@{}}
\toprule
\rowcolor{ana}\textcolor{white}{\textbf{Gereksinim}} & \textcolor{white}{\textbf{Durum}} \\
\midrule
Masal sayısı 8--12 arasında (8 masal)               & $\boxtimes$ \\
Tek HTML dosyası                                    & $\boxtimes$ \\
Stil kodları içsel (\texttt{<style>} içinde)        & $\boxtimes$ \\
Betik kodları içsel (\texttt{<script>} içinde)      & $\boxtimes$ \\
İnternetsiz açılıp çalışıyor                        & $\boxtimes$ \\
JavaScript sayfada anlamlı bir işlev görüyor        & $\boxtimes$ \\
Masal metinleri kendi cümlelerimle yazıldı          & $\boxtimes$ \\
Sayfa telefonda da düzgün görünüyor                 & $\boxtimes$ \\
\bottomrule
\end{tabular}
\end{center}

% ====================================================================
\section{Masalların Seçimi}
% ====================================================================

\begin{longtable}{@{}cp{4cm}p{3cm}p{5.4cm}@{}}
\toprule
\textbf{No} & \textbf{Masal} & \textbf{Tür} & \textbf{Seçilme nedeni} \\
\midrule
\endhead
1 & Öksüz Kız & Eğitici & Öksüz kızın zorluklarla sınandığı ve sonunda mutluluğa ulaştığı bir anlatı. \\
2 & Kurnaz Kurt & Hayvan & Köylüler ile kurdun pazarlığını konu alan kısa bir hayvan masalı. \\
3 & Tembel Çocuk ve Padişah Kızı & Eğitici & Tembel bir oğlanın çalışmayı öğrenerek değiştiği bir hikâye. \\
4 & Mitişik Kız & Fantastik & Dev anasını akıllılıkla alt eden bir kızın hikâyesi. \\
5 & Korkak Adam ve Devler & Fantastik & Korkak bir adamın devleri kurnazlıkla yendiği bir anlatı. \\
6 & Kırk Dost Masalı & Eğitici & Çıkara dayalı dostluğu konu alan, sonunda ders çıkarılan bir masal. \\
7 & Üç Haylaz Masalı & Hayvan & Çiftlikten kaçan üç hayvanın başına gelenleri anlatan, ders içeren bir masal. \\
8 & İncili Çadır & Fantastik & Yeraltı sarayı, incili çadır ve kılık değiştirme gibi masal ögeleri içeren bir anlatı. \\
\bottomrule
\end{longtable}

% ====================================================================
\section{Tasarım Kararları}
% ====================================================================

\subsection{Konuya Uygunluk}
Sayfanın ana fikri, masalları fiziksel bir kitaplık gibi sunmaktır. Her masalın kendi
renkli kapağı olan bir kitabı vardır ve kitap tıklanınca açılır. Kapak renkleri
(krem, adaçayı yeşili, lila, gül, pudra, mürdüm, turkuaz, mavi) canlı değil, yumuşak ve
mat tonlardan seçilmiştir. Masal metinleri ve kapak başlıkları, kitap baskısını
çağrıştıran serifli (Georgia) yazı tipiyle; arayüz ögeleri ise sade bir sistem yazı
tipiyle yazılmıştır. Her kitabın altında ``Anonim Türk Masalı'' ibaresi yer alır.

\subsection{Renk Paleti ve Yazı Tipi}
\begin{center}
\rowcolors{2}{acik}{white}
\begin{tabular}{@{}llp{6.4cm}@{}}
\toprule
\rowcolor{ana}\textcolor{white}{\textbf{Ad}} & \textcolor{white}{\textbf{Kod}} & \textcolor{white}{\textbf{Nerede / neden kullandınız?}} \\
\midrule
Zemin (açık gri)   & \texttt{\#fbfbfb} & Sayfa zemini; metinle yüksek karşıtlık sağlar. \\
Ana metin (siyah)  & \texttt{\#111111} & Başlıklar ve arayüz yazıları. \\
Soluk metin (gri)  & \texttt{\#777777} & Yazar satırı ve ikincil bilgiler. \\
Çizgi (açık gri)   & \texttt{\#e2e2e2} & Bölüm ayırıcı çizgiler ve kutu kenarlıkları. \\
Kapak renkleri     & \texttt{\#f4ebd0}, \texttt{\#93a193}, \texttt{\#b5adc0}, \texttt{\#d69f93}, \texttt{\#f2c6d0}, \texttt{\#5c4154}, \texttt{\#54737a}, \texttt{\#8da4c4} & Her masala bir kapak rengi verilir; kitapları birbirinden ayırır. \\
Georgia (serifli)  & -- & Masal metinleri ve kitap kapağı başlıkları. \\
Sistem yazı tipi (sans-serif) & -- & Menü, düğmeler ve alt bilgi. Sayfa internetsiz çalışacağı için harici yazı tipi kullanılmamıştır. \\
\bottomrule
\end{tabular}
\end{center}

\subsection{Okunabilirlik, Gezinme ve Duyarlı Tasarım}
\begin{itemize}
  \item \textbf{Okunabilirlik:} Okuma ekranında metin genişliği en fazla 700 piksel ile sınırlandırılmıştır. Satır aralığı 1,8, varsayılan yazı boyutu 1,1~rem'dir. Koyu metin açık zemin üzerinde yer alır. Metin iki yana yaslıdır.
  \item \textbf{Gezinme:} Kullanıcı kütüphane ekranında başlığa göre arama yapabilir veya tür listesinden süzme yapabilir. Bir kitaba tıklayınca masal açılır; ``Geri Dön'' düğmesi kütüphaneye döndürür. ``Hakkında'' düğmesi proje bilgisini açıp kapatır.
  \item \textbf{Duyarlı tasarım:} Kitaplar CSS grid ile dizilir. Ekran genişliği 768 pikselin altına inince üst menü dikey hâle gelir, kitap sütunları daralır ve okuma ekranı kenar boşlukları küçülür.
  \item \textbf{Tutarlılık:} Köşeleri yuvarlatılmış ``kapsül'' biçimli giriş kutusu, düğme ve seçim kutusu aynı kenarlık ve renklerle tasarlanmıştır. Renkler ve boşluklar CSS değişkenleri ve ortak sınıflarla yönetilmiştir.
\end{itemize}

\subsection{Ekran Görüntüleri}
% Dosyalar Overleaf projesinin ana klasöründe şu adlarla durmalı:
%   ekran-masaustu.png  ve  ekran-mobil.png
\begin{center}
  \begin{minipage}[b]{0.64\textwidth}
    \centering
    \ekrangorsel{ekran-masaustu.png}{\linewidth}\\[4pt]
    {\small Bilgisayar görünümü}
  \end{minipage}\hfill
  \begin{minipage}[b]{0.30\textwidth}
    \centering
    \ekrangorsel{ekran-mobil.png}{\linewidth}\\[4pt]
    {\small Telefon görünümü}
  \end{minipage}
  \par\smallskip
  \captionof{figure}{Sayfanın kütüphane görünümü: bilgisayar (solda) ve telefon (sağda)}
\end{center}

% ====================================================================
\section{Gerçekleştirme}
% ====================================================================

\subsection{HTML Yapısı}
Sayfa dört ana bölümden oluşur: üst bilgi (\texttt{<header>}; logo, arama kutusu ve
``Hakkında'' düğmesi), kütüphane görünümü (başlık, tür süzgeci ve kitapların dizildiği
\texttt{<main>}), okuma görünümü (araç çubuğu ve masal metni alanı) ve alt bilgi
(\texttt{<footer>}). ``Hakkında'' metni ayrı bir \texttt{<section>} içindedir ve
başlangıçta gizlidir. Masal metinleri, sayfada görünmeyen bir \texttt{<div>} içinde
\texttt{<article>} etiketleriyle tutulur. Her \texttt{<article>} etiketinde
\texttt{data-title}, \texttt{data-theme} ve \texttt{data-category} öznitelikleri bulunur.

\subsection{İçsel CSS}
Renkler \texttt{:root} içinde değişken olarak tanımlanmıştır. Kitap kapakları
\texttt{transform-style: preserve-3d} ve \texttt{rotateY(-160deg)} ile açılır.
Kitaplar \texttt{grid} ile, üst menü ve araç çubuğu \texttt{flex} ile dizilmiştir.
\texttt{@media (max-width: 768px)} kuralı telefon düzenini belirler.

\begin{lstlisting}[caption={İçsel CSS'ten önemli bir bölüm}]
#library-view {
    display: grid;
    grid-template-columns: repeat(auto-fill, minmax(220px, 1fr));
    gap: 4rem 3rem;
}

@media (max-width: 768px) {
    #library-view {
        grid-template-columns: repeat(auto-fill, minmax(130px, 1fr));
        gap: 2.5rem 1rem;
    }
}
\end{lstlisting}

\subsection{İçsel JavaScript}
Masal verisi, HTML içindeki gizli \texttt{<article>} etiketlerinde tutulur. Sayfa
yüklenince \texttt{querySelectorAll} ile bu etiketler okunur ve her biri için bir
kitap (kapak, arka kapak ve başlık bilgisi) \texttt{createElement} ile oluşturulup
kütüphaneye eklenir. Kitaba tıklanınca \texttt{is-open} sınıfı eklenir; 800 milisaniye
sonra kütüphane gizlenir, okuma görünümü açılır ve masalın metni okuma alanına
yazılır. Arama kutusundaki \texttt{input} olayı ve tür listesindeki \texttt{change}
olayı \texttt{filtreleriUygula} fonksiyonunu çalıştırır. Bu fonksiyon, başlığı aranan
kelimeyi içeren ve seçilen türe uyan kitapları gösterir, diğerlerini gizler.
\texttt{fontDegistir} fonksiyonu okuma yazısını 0,9~rem ile 1,8~rem arasında
0,15~rem adımlarla değiştirir.

\begin{lstlisting}[caption={İçsel JavaScript'ten önemli bir bölüm}]
function filtreleriUygula() {
    const arananKelime = aramaKutusu.value.toLowerCase();
    const secilenKategori = kategoriFiltresi.value;
    const tumKitaplar = document.querySelectorAll('.book-wrapper');

    tumKitaplar.forEach(kitap => {
        const kitapBasligi = kitap.getAttribute('data-title');
        const kitapKategorisi = kitap.getAttribute('data-category');

        const kelimeUyuyor = kitapBasligi.includes(arananKelime);
        const kategoriUyuyor = (secilenKategori === 'tumu' || secilenKategori === kitapKategorisi);

        if(kelimeUyuyor && kategoriUyuyor) {
            kitap.style.display = 'flex';
        } else {
            kitap.style.display = 'none';
        }
    });
}
\end{lstlisting}

\subsection{Sayfanın Özellikleri}
\begin{itemize}
  \item Üç boyutlu, açılan kitap kapağı efekti.
  \item Masal başlığına göre arama.
  \item Hayvan, eğitici ve fantastik masallara göre süzme.
  \item Okuma ekranında yazı boyutunu küçültme ve büyütme (A$-$ / A+).
  \item ``Hakkında'' bölümünü açıp kapatma ve Türkiye Kültür Portalı bağlantısı.
  \item Okuma ekranından kütüphaneye geri dönme.
  \item 768 pikselin altındaki ekranlar için duyarlı düzen.
\end{itemize}

% ====================================================================
\section{Yapay Zekâ Kullanımı}
% ====================================================================

\begin{longtable}{@{}p{2.6cm}p{3.6cm}p{4.2cm}p{3.4cm}@{}}
\toprule
\textbf{Araç} & \textbf{Ne için kullandınız?} & \textbf{İstem (prompt) özeti} & \textbf{Sayfada nerede?} \\
\midrule
\endhead
Gemini & Kod (fonksiyonların çalışması) & Sayfadaki fonksiyonlar çalıştırıldığında kodun ne yaptığını ve nasıl işlediğini açıklamasını istedim. & İçsel JavaScript bölümü (Bölüm 4.3) \\
\bottomrule
\end{longtable}

\subsection{Anladığım ve Kendim Yaptığım Kısımlar}
Gemini'yi, sayfadaki fonksiyonlar çalıştığında kodun neyi nasıl yaptığını anlamak
için kullandım. Bu açıklamalar sayesinde olayların (tıklama, \texttt{input},
\texttt{change}) hangi fonksiyonu çalıştırdığını ve verinin gizli
\texttt{<article>} etiketlerinden kitaplara nasıl dönüştüğünü sunumda anlatabilecek
düzeye getirdim.

Ayrıca koda eklenebilecek açık kaynaklı çözümleri inceledim. Fuse.js, yazım hatalarını
tolere eden, bağımlılığı olmayan ve Apache 2.0 lisanslı bir
bulanık arama kütüphanesidir\footnote{Fuse.js:
\url{https://github.com/krisk/Fuse}, erişim tarihi: 5 Kasım 2026.}.
mark.js, arama sonuçlarında anahtar kelimeleri vurgulamak için kullanılabilen, MIT
lisanslı bir kütüphanedir\footnote{mark.js:
\url{https://github.com/julkue/mark.js}, erişim tarihi: 5 Kasım 2026.}.
Bu kütüphaneler masal metni içinde arama ve eşleşen kelimeyi vurgulama gibi
geliştirmelere uygun olabilir; bu benim değerlendirmemdir. Ancak ödev yönergesi sayfanın
internetsiz çalışmasını istediği ve harici kütüphane kullanılmamasını söylediği için
bunları sayfaya eklemedim. Türkçe harflerin aramada doğru eşleşmesi için ise
JavaScript'in kendi \texttt{toLocaleLowerCase('tr')} yönteminin kullanılabileceğini
öğrendim\footnote{MDN Web Docs, \emph{String.prototype.toLocaleLowerCase()}:
\url{https://developer.mozilla.org/en-US/docs/Web/JavaScript/Reference/Global_Objects/String/toLocaleLowerCase},
erişim tarihi: 5 Kasım 2026.}.

% ====================================================================
\section{Haftalık Katkılar}
% ====================================================================

\begin{longtable}{@{}>{\raggedright\arraybackslash}p{1.7cm}%
                    >{\raggedright\arraybackslash}p{6.4cm}%
                    >{\centering\arraybackslash}p{6cm}@{}}
\toprule
\rowcolor{ana}\textcolor{white}{\textbf{Hafta}} & \textcolor{white}{\textbf{Yapılan çalışma}} & \textcolor{white}{\textbf{Görsel}} \\
\midrule
\endfirsthead
\toprule
\rowcolor{ana}\textcolor{white}{\textbf{Hafta}} & \textcolor{white}{\textbf{Yapılan çalışma}} & \textcolor{white}{\textbf{Görsel}} \\
\midrule
\endhead
\bottomrule
\endfoot

\textbf{1. hafta}\newline 07.10.2026 &
WorldWideWeb sayfasının temel HTML etiketleri (h1, p, ul vb.) ile ilk iskeleti oluşturuldu.\newline \textbf{Commit:} \href{https://drive.google.com/drive/u/0/folders/1Wlyv1icZobCmtnveznaIZm8JtNZ-Ol32}{Drive Bağlantısı}\katilimkutulari &
\haftagorsel{hafta12jpg.jpg} \\
\midrule
\textbf{2. hafta}\newline 07.10.2026 &
WorldWideWeb sayfası içsel CSS ile biçimlendirilerek renk ve tipografi ayarları yapıldı.\newline \textbf{Commit:} \href{https://drive.google.com/drive/u/0/folders/1Wlyv1icZobCmtnveznaIZm8JtNZ-Ol32}{Drive Bağlantısı}\katilimkutulari &
\haftagorsel{hafta12ajpg.jpg} \\
\midrule
\textbf{3. hafta}\newline 07.10.2026 &
Frutiger Aero temalı web arayüzü için CSS Grid/Flexbox düzenleri ve form elemanları kodlandı.\newline \textbf{Commit:} \href{https://drive.google.com/drive/u/0/folders/1u9egfAF82-hVE33o2iLil3hkjuRdCLM1}{Drive Bağlantısı}\katilimkutulari &
\haftagorsel{hafta3.jpg} \\
\midrule
\textbf{4. hafta}\newline 07.10.2026 &
HTML ve CSS ile "Shrek Sineması" bilgi kartı sayfası oluşturuldu; \texttt{<figure>} ve \texttt{<figcaption>} yapısı kullanıldı.\newline \textbf{Commit:} \href{https://drive.google.com/drive/u/0/folders/1ITVjYrJb51h5vZ_Pc2U5IV5w4n8aCY6M}{Drive Bağlantısı}\katilimkutulari &
\haftagorsel{hafta4x3.jpg} \\
\midrule
\textbf{5. hafta}\newline 21.10.2026 &
"Çiçek Eşleştirme" kart oyunu arayüzü CSS Grid ile dizayn edildi ve hover efektleri eklendi.\newline \textbf{Commit:} \href{https://drive.google.com/drive/u/0/folders/1N-IhR0fx5eB_lX_PL0x4CovzYzwQvaWV}{Drive Bağlantısı}\katilimkutulari &
\haftagorsel{hafta5.jpg} \\
\midrule
\textbf{6. hafta}\newline 21.10.2026 &
JavaScript \texttt{onclick} fonksiyonu ile "Işığı Aç / Söndür" uygulaması geliştirildi ve Bootstrap temelleri incelendi.\newline \textbf{Commit:} \href{https://drive.google.com/drive/u/0/folders/15i-x3up5RjHJ99sYIgIHuNVZrH_m3VTH}{Drive Bağlantısı}\katilimkutulari &
\haftagorsel{hafta6x2.jpg} \\
\midrule
\textbf{7. hafta}\newline [tarih] &
[Bu hafta yaptığınız çalışma.]\newline \textbf{Commit:} [mesaj veya bağlantı]\katilimkutularibos &
\haftagorsel{hafta7.png} \\
\midrule
\textbf{8. hafta}\newline [tarih] &
[Bu hafta yaptığınız çalışma.]\newline \textbf{Commit:} [mesaj veya bağlantı]\katilimkutularibos &
\haftagorsel{hafta8.png} \\
\end{longtable}

% ====================================================================
\section{Test ve Değerlendirme}
% ====================================================================

\begin{center}
\rowcolors{2}{acik}{white}
\begin{tabular}{@{}p{6cm}p{3.4cm}p{3.4cm}@{}}
\toprule
\rowcolor{ana}\textcolor{white}{\textbf{Özellik}} & \textcolor{white}{\textbf{Bilgisayar}} & \textcolor{white}{\textbf{Telefon}} \\
\midrule
Kitap kapağı açılması ve masal okuma & Çalışıyor & Çalışıyor \\
Başlığa göre arama                   & Çalışıyor & Çalışıyor \\
Türe göre süzme                      & Çalışıyor & Çalışıyor \\
Yazı boyutu (A$-$ / A+)              & Çalışıyor & Çalışıyor \\
Hakkında bölümü                      & Çalışıyor & Çalışıyor \\
\bottomrule
\end{tabular}
\end{center}

\textbf{Denenen tarayıcılar:} Firefox ve Safari.

% ====================================================================
\section{Sunum Videosu}
% ====================================================================
\begin{center}
\begin{tabular}{@{}ll@{}}
\textbf{Video süresi:} & 15:30 \quad \aciklama{12--18 dakika olmalıdır} \\
\textbf{Görünürlük:} & Herkese açık \\
\textbf{Bağlantı:} & \youtubelink \\
\end{tabular}
\end{center}

\subsection{Video Kontrol Listesi}
\aciklama{Teslimden önce hepsini kontrol edin ve kutuları $\boxtimes$ ile işaretleyin.}
\begin{itemize}[label={}]
  \item $\boxtimes$ Süre 12--18 dakika aralığında.
  \item $\boxtimes$ Sunum boyunca kameram açık, yüzüm görünüyor.
  \item $\boxtimes$ Görünürlük \textbf{Listelenmemiş} (ya da Herkese açık); \textbf{Özel} değil.
  \item $\boxtimes$ Bağlantıyı gizli pencerede açtım, video oynuyor.
  \item $\boxtimes$ YouTube videoyu işlemeyi bitirdi.
  \item $\boxtimes$ Videoyu baştan sona izledim; ses ve kamera tüm sunumda var.
  \item $\boxtimes$ Kapaktaki bağlantı ile bu bölümdeki bağlantı aynı.
\end{itemize}

% ====================================================================
\section{Sonuç ve Kendi Değerlendirmem}
% ====================================================================
Bu projede sekiz Türk masalını, tek bir HTML dosyası içinde HTML, içsel CSS ve içsel
JavaScript ile okunabilir bir web sayfasına dönüştürdüm. Masalları fiziksel bir
kitaplık gibi sunan üç boyutlu kitap kapakları, sayfanın temel tasarım fikri oldu.
JavaScript ile masal verisinin kitaplara dönüştürülmesini, başlığa göre aramayı, türe
göre süzmeyi ve okuma ekranındaki yazı boyutu ayarını öğrendim ve uyguladım. Sayfayı
bilgisayarda ve telefonda, Firefox ve Safari tarayıcılarında denedim. İleride aramayı
yalnızca başlıkla sınırlı bırakmayıp masal metni içinde de yapmak ve okuyucunun yazı
boyutu tercihini hatırlamak gibi geliştirmeler eklemek isterim.

% ====================================================================
\section{Kaynakça}
% ====================================================================
\begin{thebibliography}{9}
\bibitem{k1} T.C. Kültür ve Turizm Bakanlığı. \emph{Türkiye Kültür Portalı: Masallar}.
\url{https://www.kulturportali.gov.tr/portal/masallar-1}, erişim tarihi: 5 Kasım 2026.
\end{thebibliography}

% ====================================================================
\appendix
\section{Teslim Kontrol Listesi}
% ====================================================================
\begin{itemize}[label={}]
  \item $\boxtimes$ MEBİS'e \textbf{yalnızca bu raporun PDF'ini} yükleyeceğim. ZIP ve HTML dosyası yüklemeyeceğim.
  \item $\boxtimes$ PDF dosya adı \texttt{OgrenciNo\_AdSoyad\_vize.pdf} biçiminde.
  \item $\boxtimes$ Kapaktaki GitHub ve YouTube bağlantıları doğru; PDF'te tıklayıp denedim.
  \item $\boxtimes$ GitHub deposunda web sayfası (tek \texttt{.html}), rapor PDF'i, \texttt{.tex} kaynakları ve \texttt{README.md} var.
  \item $\boxtimes$ GitHub deposu \textbf{herkese açık} (ya da öğretim üyesine erişim verdim).
  \item $\boxtimes$ Çalışmamın son hâli teslim saatinden \textbf{önce} depoya gönderildi (sonraki commitler değerlendirilmez).
  \item $\boxtimes$ Sayfa içsel CSS ve içsel JavaScript kullanıyor, harici dosya gerektirmiyor.
  \item $\boxtimes$ Hazır taslak kullanmadım; yapay zekâ kullandıysam Bölüm 5'i doldurdum.
  \item $\boxtimes$ Çalışmanın mantığını sunumda anlatabilirim.
\end{itemize}

\end{document}